\documentclass[11pt]{article}
\usepackage[margin=1.1in]{geometry}
\usepackage{amsmath,amssymb,amsthm}
\usepackage{xcolor}
\usepackage[colorlinks=true,linkcolor=blue,citecolor=blue,urlcolor=blue]{hyperref}

\newtheorem{theorem}{Theorem}[section]
\newtheorem{proposition}[theorem]{Proposition}
\newtheorem{lemma}[theorem]{Lemma}
\newtheorem{corollary}[theorem]{Corollary}
\newtheorem{definition}[theorem]{Definition}
\newtheorem{hypothesis}[theorem]{Hypothesis}
\newtheorem{heuristic}[theorem]{Heuristic}
\theoremstyle{remark}
\newtheorem{remark}[theorem]{Remark}

\newcommand{\Z}{\mathbb{Z}}
\newcommand{\vt}[1]{v_2\!\left(#1\right)}
\newcommand{\vth}[1]{v_3\!\left(#1\right)}
\newcommand{\w}{\omega}
\newcommand{\wnext}{\w_{+}}
\newcommand{\dnext}{d_{+}}
\newcommand{\LL}{\log_2 3}

\title{Reduced coordinates for the Collatz map:\\ exact per-step laws, anchor dynamics,\\ and the limits of counting arguments for cycles}
\author{Ben Macindoe\thanks{\texttt{macindoebenjamin@gmail.com}. Developed through extensive AI-assisted research (Claude, Anthropic); the division of labor and the verification protocol are stated in Appendix~\ref{app:method}.}}
\date{v3, August 2026 $\;\cdot\;$ DOI: \href{https://doi.org/10.5281/zenodo.21730505}{10.5281/zenodo.21730505}}

\hypersetup{
  pdftitle={Reduced coordinates for the Collatz map: exact per-step laws, anchor dynamics, and the limits of counting arguments for cycles},
  pdfauthor={Ben Macindoe}
}

\begin{document}
\maketitle

\begin{abstract}
We study the Collatz dynamics in a reduced coordinate system that compresses each deterministic valuation run into a single block, producing a self-map $F$ on states $(\w,d)$ --- an odd core prime to $3$ and a depth --- whose convergence problem and cycle set are exactly those of the Collatz map. In these coordinates the local arithmetic admits exact laws, all governed by one $2$-adic quantity, the \emph{anchor} $M(\w) = -2\log\w/\log 9 \in \Z_2$: the exit valuation obeys the global law $s = 2 + \vt{d - M(\w)}$ on the lifting classes and is constant on the rest; the depth evolution closes exactly in the anchor displacement together with a stated $3$-adic absorption law; and the anchor increment obeys an exact law modulo any power of $2$, computable from graded residues of the state. A finite window of digits and stratum labels consequently decides each step in an error-free trichotomy, while a digit-budget accounting indicates that no bounded window can decide infinite horizons --- localizing the difficulty in the digit supply of the anchors. On the cycle side, a one-line elimination identity yields short rederivations of the classical exclusions at one, two and three blocks, and our main new theorem is a sharp dichotomy for counting arguments: a trim uniform in the number of blocks $p$ gives effective finiteness at every period, but its constant necessarily degrades like $(\LL)^{-p}$, and an explicit family of near-counterexamples (\emph{staircases}: divergent-orbit profiles bent into loops) shows counting cannot do substantially better, so uniform cycle exclusion requires arithmetic (divisibility) input, not sharper counting. Finally we state the equidistribution hypothesis implicit in the classical heuristics as a precise conjecture --- that the block letters of uniformly sampled large starts carry, at every finite block length, the frequencies of an exactly computable Bernoulli law --- prove its consequences conditionally, exhibit an unconditional base case inside the digit budget, and report a calibration campaign whose four apparent anomalies all dissolved under controls.
\end{abstract}

\subsection*{Version note}
v1, July 2026 (original publication). v2, July 2026 (v2-specific Zenodo DOI: 10.5281/zenodo.21421120): the subtitle of the \texttt{merle} citation restored, and a note added evidencing the sharpness hedge of Theorem~\ref{thm:staircase}, prompted by correspondence with Eric Merle. v3, August 2026 (drafted; the version-specific DOI on the title page is reserved and this version is not yet published): after external review, each definition, statement and scope word is brought back into line with the project record; the sharpness assessment of Section~\ref{sec:cycles} is reported as proved there; and Section~\ref{sec:aeh} is restated around one object and one hypothesis in ensemble form. No numbered theorem's claim is strengthened, weakened, or renumbered in either revision, and nothing new is proved in this paper; v3 reports two results established in the project record and not reproduced here. The repair history, item by item, is in this version's release description.
\subsection*{Author's note}

\emph{Dear reader --- I'm no mathematician, and I learned much of this paper's terminology during the process of formalizing it ($p$-adic valuations were completely new to me as a formal concept). It all began when I got hooked on a strange recursive pattern in the prime factors of a Collatz chain, and that idea pushed me into conversation with AI. We cleaned my idea up into the states $(\w,d)$, with the primary goal of factoring as much deterministic behaviour as possible out of the question, in the hope of focusing attention where it matters. That simplification is what enables the paper's main contributions: exact laws for every single step; a uniform trim that caps what size-counting can achieve against cycles; and an exact hypothesis pointing the remaining question toward more fundamental ``fair-coin'' mathematics. Because this work is AI-assisted and because I am not formally trained in the relevant mathematics, Appendix~\ref{app:method} describes the verification protocol used throughout the project. The complete research record is public at \url{https://github.com/macindoe/collatz}. --- Ben}

\section{Introduction}

Let $T(x) = (3x+1)/2^{\vt{3x+1}}$ denote the odd-to-odd Collatz map on positive odd integers. The Collatz conjecture asserts that every $T$-orbit reaches $1$. This paper develops a coordinate system in which the local arithmetic of $T$ becomes exactly computable, and uses it for three purposes: to prove exact per-step laws for the induced dynamics; to determine how far counting arguments can go toward cycle exclusion (answer: exponentially far in the number of blocks, and provably no farther); and to state precisely the statistical hypothesis that the classical heuristics leave implicit.

The idea is to stop watching individual steps. Writing $x + 1 = 2^m u$ with $u$ odd, the next $m-1$ odd steps of $T$ are deterministic --- each converts one factor of $2$ in $x+1$ into a factor of $3$ --- and terminate in a single cascade (Proposition~\ref{prop:block}). Compressing this \emph{block} and splitting $u = 3^a\w$ with $3\nmid\w$ yields a state $(\w, d)$ with $d = m + a$, and the block-to-block dynamics induce a self-map $F$ on such states. The compression loses nothing the conjecture depends on (Theorem~\ref{thm:equiv}).

\paragraph{Contributions.} (i) A reduced self-map $F$ on states $(\w,d)$, with an in-paper proof of well-definedness, whose convergence problem and cycle set are exactly equivalent to those of $T$ (Section~\ref{sec:reduced}). (ii) A global anchor law for the exit valuation $s$, derived from one squaring lemma (Section~\ref{sec:anchor}). (iii) Exact residue-class formulas for the depth evolution, including an explicit $3$-adic absorption law (Theorem~\ref{thm:depth}). (iv) An exact low-order law for the anchor increment and an error-free trichotomy deciding the next step from a finite residue window together with the step's stratum labels, with a digit-budget accounting (Theorems~\ref{thm:deltaM}--\ref{thm:onestep}, Heuristic~\ref{prop:budget}). (v) A cycle elimination identity, short rederivations of the classical small-cycle exclusions, and --- the main new theorems --- a uniform trim with effective finiteness at every period together with a sharpness family showing the exponential loss is intrinsic to size-counting (Section~\ref{sec:cycles}). (vi) The equidistribution hypothesis stated exactly, its consequences proved conditionally, and a calibration record (Section~\ref{sec:aeh}).

Section~\ref{sec:anchor} is the technical heart: every quantity above obeys an exact law in the anchor displacement $d - M(\w)$, a finite window of digits and stratum labels then decides the next step in a trichotomy that never errs, and an elementary accounting argument shows each decision \emph{consumes} digits irreversibly. In these coordinates the difficulty of the problem is not diffuse: it is the question of where the digits of the anchors come from. Section~\ref{sec:cycles} tests the one regime where finite data can beat unbounded depth --- cycles --- and, since the staircase is precisely a divergent-orbit profile closed into a loop, the cycle half and the divergence half of the residual difficulty are, in a concrete sense, one object. Section~\ref{sec:aeh} formalizes the statistical half: the folklore ``fair-coin'' heuristics (Terras \cite{terras}) become an exact statement about the block letters of typical orbits, equivalently about an explicitly computable window law $\pi_{k,D}$, whose ledger consequences are unconditional theorems \emph{about} $\pi_{k,D}$.

\subsection*{Related work and provenance}
Block-like decompositions and stochastic models are classical (Terras \cite{terras}; see Lagarias \cite{lagarias} for the literature), and the $2$-adic conjugacy of the Collatz map to the shift underlies all Haar-genericity heuristics. The unconditional density line descending from Terras is long and is not what Section~\ref{sec:aeh} adds to. Terras \cite{terras} and Everett \cite{everett} place descent below the start at natural density one; Korec \cite{korec} lowers the target to $y^{c}$ for every $c > \log_4 3 = 0.7924\ldots$, in natural density and within $O(\log y)$ steps; Inselmann \cite{inselmann} carries a two-sided trajectory law --- the orbit tracked to within $m^{\pm\varepsilon}$ of its classical drift, simultaneously at every time up to $(\log_2\tfrac43)^{-1}\log_2 m$ Syracuse steps, the full descent horizon in that unit --- at natural density one, and so reaches $m^{\varepsilon}$ for every fixed $\varepsilon>0$; and Tao \cite{tao} reaches every divergent threshold, in logarithmic density. All of these horizons are counted in steps, and their reading in the reduced blocks of Section~\ref{sec:aeh} is not free: it divides by a mean number of steps per block which is a theorem of the cylinder count inside the digit budget and, past it, neither a theorem nor a consequence of Hypothesis~\ref{hyp:aeh}. What Hypothesis~\ref{hyp:aeh} asserts is not descent but the letter statistics themselves --- equivalently the exact window law $\pi_{k,D}$ --- at horizons past the digit budget of Heuristic~\ref{prop:budget}. Tao's paper carries a footnote to two further $3$-adic studies, and neither is the object appearing here. Wirsching \cite{wirsching} completes the root of the predecessor tree $3$-adically and studies the inverse map as a Markov chain on $[0,1] \times \Z_3^{\times}$ in order to \emph{count} predecessors; his inverse branches therefore carry equal weight and his limiting $3$-adic law is Haar, whereas the law appearing in Section~\ref{sec:aeh} weights each branch by the probability $2^{-a}$ that an orbit takes it and is consequently not uniform --- $(\tfrac23,\tfrac13)$ at level one where Haar gives $(\tfrac12,\tfrac12)$. The two programmes are nonetheless siblings in shape: what Wirsching needs is that a Collatz-generated object be unbiased in $3$-adic coordinates, and Hypothesis~\ref{hyp:aeh} is that one is unbiased in $2$-adic coordinates. Thomas \cite{thomas} proves that the integers whose orbit reaches $1$ with all six frequencies over the units modulo $18$ below $\tfrac16 + 0.0215$ have counting function at most $N^{0.9999}$; he constructs no limiting measure, and the law recorded here implies his conclusion, its level-one form already forcing some class to frequency at least $\tfrac29$. Two-adic and Diophantine methods dominate the cycle literature \cite{steiner,sdw,hercher,eliahou}; effective linear forms in logarithms \cite{yu,bl,rhin} are the engine, and the verified range of the conjecture stands at $2^{71}$ \cite{barina}. Two contemporary human--AI efforts are close in spirit: a documented exploration \cite{llmcollatz} develops a burst--gap decomposition and an orbit equidistribution conjecture, and a Lean~4 formalization \cite{merle} proves cycle exclusion conditionally on Baker-type and verification hypotheses, including an impossibility lemma for measure-based uniform bounds kindred to our Theorem~\ref{thm:staircase}; neither contains the anchor machinery, exact per-step laws, or the counting-limit dichotomy developed here. A bachelor's thesis of Rackl (Klagenfurt, 2021) applies valuation methods to Syracuse cycles via the classical cycle equation; it likewise contains none of the present apparatus. The division of labor in the present work is stated precisely, because it matters for both credit and blame. The reduced coordinate system itself --- the odd core (``seed''), the block cascade, the states $(\w,d)$, and the driving question of Section~\ref{sec:anchor} --- is the human author's invention, originating in a hand-drawn note of 13~March~2023 that predates the collaboration. The analytical superstructure erected on it --- the anchor formulation, the cycle apparatus, the equidistribution formalization, and the proofs --- was developed in AI-assisted research under the author's direction, drawing on the classical literature's concepts; any overlap with contemporary work lies in that layer, not in the coordinate system. Appendix~\ref{app:method} states the responsibility and verification protocol.

\section{The reduced system}\label{sec:reduced}

\begin{definition}\label{def:reduced}
A \emph{reduced state} is a pair $(\w, d)$ with $\w$ odd and positive, $3\nmid\w$, $d \ge 1$. Set
\[ A = 3^d\w - 1,\qquad s = \vt{A},\qquad x_{\mathrm{exit}} = A/2^s,\qquad C = A + 2^s, \]
\[ \sigma = \vt{C},\qquad a_+ = \vth{C},\qquad m_+ = \sigma - s,\qquad \wnext = \frac{C}{2^{\sigma}3^{a_+}},\qquad \dnext = m_+ + a_+, \]
and define the \emph{reduced map} $F(\w,d) = (\wnext, \dnext)$. The \emph{projection} $R$ sends an odd $x$ with $x + 1 = 2^m 3^a \w$ ($3\nmid\w$) to $R(x) = (\w, m + a)$.
\end{definition}

\begin{proposition}[block structure and well-definedness]\label{prop:block}
Let $x$ be odd with $x + 1 = 2^m u$, $u = 3^a\w$ odd, $3 \nmid \w$, and let $(\w,d) = R(x)$, $d = m+a$. Then:
\begin{enumerate}
\item[(i)] the next $m - 1$ iterates of $T$ are $x_j$ with $x_j + 1 = 2^{m-j}\,3^j u$ for $j = 0,\dots,m-1$;
\item[(ii)] the following iterate is $T(x_{m-1}) = A/2^s = x_{\mathrm{exit}}(\w,d)$, which depends only on $(\w,d)$;
\item[(iii)] consequently $R(T^{m}(x)) = F(R(x))$ for every representative $x$ of $(\w,d)$, the representatives being exactly $x = 2^{d-a}3^a\w - 1$, $0 \le a \le d-1$; and the sequence of reduced states visited by any $T$-orbit is the $F$-orbit of its projection.
\end{enumerate}
\end{proposition}

\begin{proof}
(i) If $y + 1 = 2^k v$ with $k \ge 2$, $v$ odd, then $3y + 1 = 3(y+1) - 2 = 2\left(3\cdot 2^{k-1}v - 1\right)$, an odd multiple of $2$, so $T(y) = 3\cdot2^{k-1}v - 1$ and $T(y) + 1 = 2^{k-1}(3v)$: one factor of $2$ is exchanged for one factor of $3$. Induction from $y = x$ gives (i). (ii) At $j = m-1$ we have $x_{m-1} + 1 = 2\cdot 3^{m-1}u$, so $3x_{m-1} + 1 = 3(x_{m-1}+1) - 2 = 2(3^m u - 1) = 2A$ since $3^m u = 3^{m+a}\w = 3^d \w$. Thus $T(x_{m-1}) = 2A/2^{1+s} = A/2^s$, and $A$ depends only on $(\w,d)$. (iii) $x_{\mathrm{exit}} + 1 = (A + 2^s)/2^s = C/2^s$, so $R(x_{\mathrm{exit}}) = (\wnext, (\sigma - s) + a_+) = F(\w,d)$ by Definition~\ref{def:reduced}; the representative list is immediate from $x + 1 = 2^{d-a}3^a\w$.
\end{proof}

\begin{theorem}[convergence translation]\label{thm:equiv}
$(1,1)$ is a fixed point of $F$ whose $R$-fiber is $\{1\}$. The Collatz conjecture holds if and only if every $F$-orbit from a valid state reaches $(1,1)$, and $R$ induces a bijection between nontrivial $T$-cycles and $F$-cycles other than $\{(1,1)\}$.
\end{theorem}

\begin{proof}
For $(1,1)$: $A = 2$, $s = 1$, $x_{\mathrm{exit}} = 1$, $C = 4$, so $F(1,1) = (1,1)$; the fiber is $\{1\}$ since the representatives are $x = 2\cdot 1 - 1 = 1$ only. If the conjecture holds and $(\w,d)$ is valid, its representative $x = 2^d\w - 1$ reaches $1$, so by Proposition~\ref{prop:block}(iii) the $F$-orbit of $(\w,d)$ reaches $R(1) = (1,1)$. Conversely if every $F$-orbit reaches $(1,1)$, then any odd $x$ has a $T$-orbit visiting the fiber of $(1,1)$, i.e.\ the value $1$. For cycles: a nontrivial $T$-cycle projects to a periodic $F$-orbit avoiding $(1,1)$ (its values avoid $1$); conversely an $F$-cycle's block entries $e_t = x_{\mathrm{exit}}(S_{t-1})$ are deterministic in the states and joined by genuine $T$-iteration, so $e_0$ is $T$-periodic and nontrivial; the two constructions are mutually inverse.
\end{proof}

\section{Anchor dynamics: the exact per-step laws}\label{sec:anchor}

For $u \equiv 1 \pmod 8$ the equation $9^{t} = u^{-1}$ has a unique solution $t \in \Z_2$, since $9$ topologically generates $1 + 8\Z_2$; write $N(u)$ for it, so $N(u) = -\log u/\log 9$ in the $2$-adic logarithm. $N$ is an isomorphism $(1+8\Z_2,\times)\to(\Z_2,+)$, and
\begin{equation}\label{eq:isometry}
\vt{9^t - 1} = 3 + \vt{t}\qquad (t \in \Z_2\setminus\{0\}),
\end{equation}
the \emph{isometry}, since $9 = 1 + 8$ and $\vt{\log(1+8z)} = 3 + \vt{z}$-type estimates hold with equality here.

\begin{definition}[anchor]
For every odd $\w$ set $M(\w) = N(\w^2) = -2\log\w/\log 9 \in \Z_2$, well defined since $\w^2\equiv 1 \pmod 8$. $M(\w)$ is even iff $\w\equiv\pm1\pmod 8$.
\end{definition}

\begin{lemma}[squaring lemma]\label{lem:squaring}
Let $\w$ be odd, $d \ge 1$, and $Y = 3^d\w$. If $Y \equiv 1 \pmod 8$, then
\[ \vt{Y - 1} = 2 + \vt{d - M(\w)}. \]
\end{lemma}

\begin{proof}
$Y^2 = 9^d\w^2 = 9^{d}\,9^{-M(\w)}\cdot\bigl(9^{M(\w)}\w^2\bigr) = 9^{\,d - M(\w)}$, since $9^{M(\w)} = \w^{-2}$ by definition of $N$. By \eqref{eq:isometry}, $\vt{Y^2 - 1} = 3 + \vt{d - M(\w)}$ (nonzero since $Y^2 \neq 1$ for a valid state). Since $Y \equiv 1 \pmod 8$ we have $\vt{Y+1} = 1$, and $\vt{Y-1} = \vt{Y^2-1} - \vt{Y+1}$ gives the claim.
\end{proof}

\begin{theorem}[global valuation law]\label{thm:vlaw}
Call $(\w,d)$ \emph{lifting} if $3^d\w \equiv 1 \pmod 8$, i.e.\ $\w\equiv1\ (8)$ with $d$ even or $\w\equiv3\ (8)$ with $d$ odd. On the lifting classes
\[ s = \vt{3^d\w - 1} = 2 + \vt{d - M(\w)} \;\ge\; 3, \]
while on the six non-lifting residue-parity classes $s \in \{1,2\}$ is the constant determined by $(\w \bmod 8,\, d \bmod 2)$: $s = 1$ on $(1,\mathrm{odd}), (5,\mathrm{odd}), (7,\mathrm{even}), (3,\mathrm{even})$ and $s = 2$ on $(5,\mathrm{even}), (7,\mathrm{odd})$. Unconditionally $s \le C(\w)\,(1 + \log d)^2$ by $p$-adic Baker theory \cite{yu,bl}.
\end{theorem}

\begin{proof}
The lifting case is Lemma~\ref{lem:squaring}; the parity claim ($d \equiv M(\w) \bmod 2$ on lifting classes) follows from $\vt{d - M(\w)} \ge 1$ there. Off-lifting, $3^d\w \bmod 8 \in \{3,5,7\}$ gives $\vt{3^d\w - 1} \in \{1,2\}$ by inspection of residues mod $8$. The Baker bound applies to $s = \vt{9^n\tilde\w - 1}$, a linear form in two $2$-adic logarithms.
\end{proof}

\begin{lemma}[$3$-adic absorption law]\label{lem:absorption}
Let $h(s) = \vth{2^s - 1}$, so $h(s) = 0$ for $s$ odd and $h(s) = 1 + \vth{s}$ for $s$ even. Then
\[ a_+ = \vth{C} = \begin{cases} 0, & s \text{ odd},\\[2pt] \min(d,\, h(s)), & s \text{ even},\ d \ne h(s),\\[2pt] d + \vth{\w + (2^s-1)3^{-d}}, & s \text{ even},\ d = h(s). \end{cases} \]
In particular the next state gains a factor of $3$ exactly when $s$ is even.
\end{lemma}

\begin{proof}
$C = 3^d\w + (2^s - 1)$ with $\vth{3^d\w} = d$ exactly and $\vth{2^s-1} = h(s)$; $h(s) = 0$ iff $s$ odd since $2^s \equiv 2 \pmod 3$ then, and for even $s$, lifting-the-exponent gives $\vth{4^{s/2}-1} = \vth{3} + \vth{s/2} = 1 + \vth{s}$. The ultrametric equality $\vth{X+Y} = \min(\vth X,\vth Y)$ holds when the valuations differ; when $d = h(s)$, factor $3^d$ and evaluate the remaining unit sum.
\end{proof}

\begin{theorem}[depth closure]\label{thm:depth}
On every residue class the next depth is $\dnext = m_+ + a_+$ with $a_+$ given exactly by Lemma~\ref{lem:absorption} and $m_+$ by the \emph{entry-depth law}: writing $\varepsilon$ for the anchor displacement ($\varepsilon = n - N(\w)$ on the lifting class $d = 2n$, $\varepsilon = n - N(3\w)$ on $d = 2n+1$) and $\delta_s = N(1 - 2^s)$ (which satisfies $\vt{\delta_s} = s - 3$ for $s \ge 3$),
\[ m_+ = 3 - s + \vt{\varepsilon + \delta_s} \quad\text{(lifting)},\qquad
   m_+ = \begin{cases} 1 + \vt{d - M(\w)}, & s = 1,\\ \vt{(d-1) - M(\w)}, & s = 2. \end{cases} \]
The forced carry $\vt{\delta_s} = s-3 = \vt{\varepsilon}$ yields $m_+ \ge 1$ structurally.
\end{theorem}

\begin{proof}
All cases reduce to the \emph{target-shift lemma}: for $c \in 1 + 8\Z_2$ and $d = 2n$ with $\w \equiv 1\ (8)$,
\[ \vt{3^{d}\w - c} \;=\; 3 + \vt{n - N(\w) + N(c)}, \]
proved exactly as Lemma~\ref{lem:squaring} after writing $3^d\w - c = 9^{-N(c)}\bigl(9^{\,n - N(\w) + N(c)} - 1\bigr)$. Since $m_+ = \vt{x_{\mathrm{exit}}+1}$ and $2^s(x_{\mathrm{exit}}+1) = C = 3^d\w - (1 - 2^s)$, the lifting case is the shift $c = 1 - 2^s$ (which lies in $1+8\Z_2$ for $s \ge 3$), with $\vt{\delta_s} = s - 3$ from the isometry; the odd lifting component reduces to the even one via the companion parameter $3\w$. For $s = 1$: the shifted target is $-1$, and $\vt{X+1} = \vt{X^2-1} - \vt{X-1}$ with $X = 3^d\w$, $X^2 = 9^{\,d - M(\w)}$, gives the formula; For $s = 2$: here $C = 3^d\w + 3 = 3(X' + 1)$ with $X' = 3^{d-1}\w$, so $m_+ = \vt{C} - 2 = \vt{X'+1} - 2$. The $s=2$ classes force $3^d\w \equiv 5 \pmod 8$, hence $X' \equiv 3^{-1}\cdot 5 \equiv 7 \pmod 8$, so $\vt{X'-1} = 1$; and $X'^2 = 9^{\,(d-1) - M(\w)}$ as in Lemma~\ref{lem:squaring}, so $\vt{X'+1} = \vt{X'^2-1} - \vt{X'-1} = \bigl(3 + \vt{(d-1) - M(\w)}\bigr) - 1$, giving $m_+ = \vt{(d-1) - M(\w)}$ as claimed.
\end{proof}

\begin{remark}[verification notes]\label{rem:verify1}
Each law above was checked by independently written code, with zero discrepancies: the valuation law on $8{,}000$ random states across both lifting components, the entry-depth law exhaustively on all $62{,}937$ valid states with $\w < 3000$, $d < 64$, and the absorption law on $327{,}980$ states ($983{,}954$ comparisons against both this Lemma and its form in the record), with the boundary branch $d = h(s)$, on which $a_+$ is a $3$-adic function of $\w$, and the shallow branch $d < h(s)$ constructed rather than sampled. The records are \texttt{stage3.md} \S11.8.6.3 and that page's status header; the code is \texttt{experiments/absorption\_law.py}.
\end{remark}

\begin{theorem}[low-order anchor-increment law]\label{thm:deltaM}
Let $\Delta M = M(\wnext) - M(\w)$. For every $k \ge 1$, $\Delta M \bmod 2^k$ is determined by, and explicitly computable from,
\[ \w \bmod 2^{\sigma + k + 2},\qquad d \bmod 2^{\sigma + k},\qquad a_+ \bmod 2^k, \]
and the stratum data $(s, \sigma)$ are determined by the same residues. This last clause does not extend to $a_+$: by Lemma~\ref{lem:absorption} the absorption is a $3$-adic function of $\w$, and it enters the display only through its residue $a_+ \bmod 2^k$.
\end{theorem}

\begin{proof}
Three digit-determinacy facts chain together. (a) For $u \equiv 1\ (8)$, $N(u) \bmod 2^k$ depends only on $u \bmod 2^{k+3}$, since $9$ has order $2^k$ mod $2^{k+3}$ and $9^t \equiv u^{-1}$ determines $t$ mod $2^k$. (b) $C \bmod 2^{q}$ is determined by $\w \bmod 2^q$, $d \bmod 2^{q-2}$, $s$, since $3$ has order $2^{q-2}$ mod $2^q$. (c) $\wnext \bmod 2^r$ is determined by $C \bmod 2^{\sigma + r}$, $\sigma$, and $a_+ \bmod 2^{r-2}$, by exact division and the order of $3$ mod $2^r$. Now $\Delta M = N\bigl((\wnext/\w)^2\bigr)$; by (a) this needs $(\wnext/\w)^2 \bmod 2^{k+3}$, hence $\wnext, \w \bmod 2^{k+2}$; by (c) with $r = k+2$ this needs $C \bmod 2^{\sigma+k+2}$ and $a_+ \bmod 2^k$; by (b) this needs the stated residues of $\w$ and $d$. Finally $s \le \sigma$ and $\sigma$ are visible in $A \bmod 2^{\sigma+1}$, $C \bmod 2^{\sigma+1}$.
\end{proof}

\begin{theorem}[one-step trichotomy]\label{thm:onestep}
The \emph{depth-$k$ window} of a valid state $(\w,d)$ consists of the residues of Theorem~\ref{thm:deltaM} together with the stratum labels $(s, \sigma, a_+)$. From the depth-$k$ window alone, one can determine: the exact next depth $\dnext$ and next residue-parity class; the exact value $s_+ \in \{1,2\}$ if the next class is non-lifting; and, if lifting, the truncation of the next displacement mod $2^k$ --- yielding $s_+$ exactly when nonzero, and the correct report ``$s_+ \ge k+2$'' when zero. The window never outputs a wrong value; under the law $\pi_{k,D}$ of Section~\ref{sec:aeh} the undecided rate is $\approx 2^{-(k+1)}$.
\end{theorem}

\begin{proof}
By facts (b), (c) above the window determines $\wnext \bmod 2^{k+2}$, hence the next class; $\dnext = (\sigma - s) + a_+$ is exact from the stratum data. Non-lifting $s_+$ is the class constant (Theorem~\ref{thm:vlaw}). Lifting: $M(\wnext) \bmod 2^k$ follows from fact (a) applied to $\wnext^2$, so $(\dnext - M(\wnext)) \bmod 2^k$ is computable; if nonzero its valuation is visible and Theorem~\ref{thm:vlaw} gives $s_+$; if zero then $\vt{\dnext - M(\wnext)} \ge k$, i.e.\ $s_+ \ge k + 2$.
\end{proof}

\begin{heuristic}[digit budget]\label{prop:budget}
Deciding one step at depth $k$ consumes the state's $2$-adic data to depth $\sigma + k + 2$ (empirical mean $\sigma \approx 4.0$ along orbits), and nothing regenerates it: the anchors of later states are functions of the full initial state, not of any bounded window. A window of $W$ digits therefore supports $O(W/(k+6))$ decided steps. This strongly suggests --- and we treat it as the organizing heuristic, not a formalized theorem --- that bounded-window determinism cannot decide unbounded horizons without additional statistical or rigidity input.
\end{heuristic}

\begin{remark}[verification note]
The trichotomy was run along real orbits for $k \in \{4,8\}$: $21{,}296$ steps, $21{,}000$ decided with zero errors, $296$ flagged deep with zero violations of the bound $s_+ \ge k+2$. Code: \texttt{experiments/one\_step\_propagation.py}; the low-order law of Theorem~\ref{thm:deltaM} is checked by \texttt{experiments/anchor\_increment.py}.
\end{remark}

Heuristic~\ref{prop:budget} is the paper's organizing negative observation: it splits the residual difficulty into a statistical question for typical starting values (Section~\ref{sec:aeh}) and a rigidity question for cycles (Section~\ref{sec:cycles}); we know of no residual content outside those two.

\section{Cycles: rederivations, a uniform trim, and its sharpness}\label{sec:cycles}

Fix a hypothetical $p$-block cycle of $F$ with entry depths $m_t \ge 1$, exit valuations $s_t \ge 1$ ($t \in \Z/p\Z$), and set $n = \sum m_t$, $K = \sum s_t + n$, $q = 2^K - 3^n$.

\begin{proposition}[elimination identity]\label{prop:elim}
$q > 0$, and for every rotation $r$,
\[ \w_r\, 3^{a_{r-1}}\, q \;=\; R_r \;:=\; \sum_{t=0}^{p-1} 3^{M_t}\, 2^{S_t}\,(2^{s_t}-1), \]
where, in rotation order starting at $r$, $M_t = \sum_{j>t} m_j$ and $S_t = \sum_{j<t}\sigma_j$, with $\sigma_j = s_j + m_{j+1}$, indices cyclic (the step's $\sigma = \vt{C}$, $m_+ = \sigma - s$, of Definition~\ref{def:reduced}; the shift is essential, $m_{j+1}$ and not $m_j$; and $M_t$, a partial sum of entry depths, is unrelated to the anchor $M(\w)$ of Section~\ref{sec:anchor} --- this is the project record's own notation, \texttt{cycles.md} \S12.6.1, kept so that paper and record read together). Hence $q \mid R_r$ and $q \le \min_r R_r$. For the trivial cycle, $R = 4^p - 3^p = q$ identically.
\end{proposition}

\begin{proof}
The step identity $2^{\sigma_t}3^{a_t}\w_{t+1} = 3^{d_t}\w_t + (2^{s_t}-1)$ (Definition~\ref{def:reduced}) unrolls $p$ times; imposing closure $\w_p = \w_0$, substituting $d_j = m_j + a_{j-1}$, and cancelling the common power of $3$ yields the display. Positivity: the right side is positive.
\end{proof}

\begin{corollary}[size balance]\label{cor:size}
$2^K = 3^n \prod_t (1 + \varepsilon_t)$ with $\varepsilon_t = (2^{s_t}-1)/(3^{d_t}\w_t)$; if every cycle element exceeds $X$ then $0 < K\log 2 - n\log 3 < p/X$. This is the classical Baker-side engine \cite{eliahou}, obtained in four lines.
\end{corollary}

\begin{lemma}[ceiling]\label{lem:ceiling}
If $K > \lceil n\LL \rceil$ then $\prod(1+\varepsilon_t) \ge 2$, so some $\varepsilon_t \ge 2^{1/p}-1 \ge (\ln 2)/p$ and some block exit is $< p/\ln 2$. Consequently $K = \lceil n \LL\rceil$ for every nontrivial cycle with $p \le \ln 2 \cdot X_0$, $X_0$ the verified range \cite{barina} --- unconditionally for $p\le5$ by inspecting the odd values $\le 7$.
\end{lemma}

\begin{theorem}[periods $1$--$3$; rederivations]\label{thm:smallp}
$F$ has no nontrivial cycles of $1$, $2$, or $3$ blocks. These statements are contained in the classical results \cite{steiner,sdw,hercher}; the derivations, not the theorems, are the contribution.
\end{theorem}

\begin{proof}[Proof outline]
$p=1$: the fixed-point equation is $\w\,3^a(2^{s+m}-3^m) = 2^s-1$; the sign and size constraints pin $2^s$ into a window of endpoint ratio $(1-3^{-m})/(1-2^{-m}) < 2$, leaving at most one candidate $s$ per $m$; exact search kills $m \le 20{,}000$ and Steiner's theorem \cite{steiner} covers all $m$ via the block--circuit correspondence. $p=2,3$: budget-trim lemmas (casework over rotations' dominant terms, using $\sum s_t = K - n$) force $0 < K\log2 - n\log3 < 2^{6-n/5}$ resp.\ $2^{3-0.115n}$; exact searches to $n = 20{,}000$ find no solutions; any effective irrationality measure for $\LL$ \cite{rhin} excludes the tail. Full details are in the project record: the three classifications at \texttt{cycles.md} \S12.2.3, \S12.5.3 and \S12.7.5, the two trim lemmas at \S12.5.2 and \S12.7.4, and the ceiling lemma that supplies $K = \lceil n\LL\rceil$ at \S12.6.2. The searches are \texttt{experiments/period1\_cycles.py}, \texttt{experiments/period2\_cycles.py} and \texttt{experiments/period3\_cycles.py}.
\end{proof}

\begin{theorem}[uniform trim]\label{thm:uniform}
Every nontrivial $p$-block cycle satisfies, with $\gamma := K - \log_2 q$,
\[ \gamma + \log_2 p \;>\; \frac{(\LL - 1)\, n}{(\LL)^{\,p} - 1}. \]
Consequently the candidate set at every period is finite and explicitly bounded. Combined with Rhin's effective bound \cite{rhin},
\[ K\log 2 - n\log 3 \;>\; \exp\bigl(-13.3\,(0.46057 + \log n)\bigr), \]
every nontrivial $p$-block cycle has $n \le n_0(p)$, the unique solution in $n$ of
\[ \frac{(\LL - 1)\,n}{(\LL)^{\,p} - 1} \;=\; \log_2 p \;+\; \frac{p + 13.3\,(0.46057 + \log n)}{\log 2}, \]
so that $n_0(p) = O(p\,(\LL)^p)$.
\end{theorem}

\begin{proof}
For each rotation, $q \le R_r < p\,2^{\max_t e_t}$, where the exact exponent identity gives $e_t - K = (\LL-1)M_t - m_r - \sum_{j>t}s_j$. The resulting conditions $m_r < \gamma + \log_2 p + \Phi_r$ close into the cyclic recursion $\Phi_{r} = \max(0, \Phi_{r-1} + w_{r-1})$ on arc weights $w_j = (\LL-1)m_j - s_j$, whose cyclic total $(\LL-1)n - \sum s_t$ is negative because $2^K > 3^n$ forces $\sum s_t > (\LL - 1)n$. A max-plus recursion with negative cyclic total must vanish somewhere (otherwise following it once around the cycle would add the total to itself); fix such a reset index $r_0$, so $\Phi_{r_0} = 0$. Walking forward, $\Phi_{r_0+j} \le (\LL-1)T_{j-1}$ where $T_j$ is the partial sum $m_{r_0} + \cdots + m_{r_0+j}$ (drop the $-s$ terms), so the block bound gives $m_{r_0+j} < (\gamma + \log_2 p) + (\LL-1)T_{j-1}$, i.e.\ $T_j < \LL\,T_{j-1} + (\gamma + \log_2 p)$. Iterating $p$ times from $T_{-1} = 0$ yields $n = T_{p-1} < (\gamma + \log_2 p)\,\bigl(\LL^{\,p} - 1\bigr)/(\LL - 1)$, which is the display.
\end{proof}

\begin{theorem}[sharpness: the staircase]\label{thm:staircase}
No trim uniform in $p$ can extend the small-period constants: there exist configurations satisfying every rotation's exact size condition $q \le R_r$ whose $\gamma$ falls far below any polynomial-in-$p$ extension of the constants of periods $2$--$3$. Explicitly, at $p=7$, $n=94$, the staircase $m = (4,7,9,15,23,35,1)$ with $s = (1,\dots,1,\,S-6)$ passes all seven size tests with $\gamma = 6.74$, where the period-$3$ constant would demand $\gamma > 8.9$; $84$ further verified instances occur at $p=6$ alone. All such configurations fail the divisibility conditions $q \mid R_r$. Uniform cycle exclusion therefore requires the divisibility system --- equivalently, rigidity of the closed anchor walk $\sum_t \Delta M_t = 0$.
\end{theorem}

\paragraph{Sharpness evidence and assessment.} The construction behind these witnesses --- block depths growing geometrically at ratio $\approx\LL$ with unit exit valuations, closed by a single block of unit depth and maximal exit valuation --- tracks the extremal configuration of the max-plus recursion in the proof of Theorem~\ref{thm:uniform}. Supported by the verified instances, though not proved here for all $p$, we assess that it passes all size conditions with $\gamma = O(\log p)$ for every $p$. This assessment is evidence for Theorem~\ref{thm:staircase} and no part of it: the theorem is closed at $p = 6, 7$ by the witnesses exhibited above. The assessment has since been proved in the project record, at a scope and a shape stronger than assessed here; the status paragraph below states what was proved, where, and at what scope.

\begin{remark}\label{rem:staircase}
The staircase profile --- geometric depth growth with shallow exits --- is exactly the growth regime of divergent orbits. Size analysis cannot forbid it as a cycle for the same reason drift analysis cannot forbid divergence. The two halves of the residual difficulty meet in one configuration.
\end{remark}

\paragraph{Status of the assessment (August 2026).} Prompted by correspondence with Eric Merle concerning this theorem's sharpness hedge (his related formal work is cited in \cite{merle}), the assessment above has since been \emph{proved} in the project record, at a scope and a shape stronger than it claims, and is not reproduced here; Theorem~\ref{thm:staircase} stands above exactly as written. A passing size-condition witness is constructed at every period, and the whole Diophantine input is one exact number: $8 - 5\LL = 0.0751874964\ldots$ is positive and below the target arc's length, so any $66$ consecutive integers contain an admissible exponent $n$, and the scale window at period $p$ holds $0.05\,(\LL)^{p}$ integers --- $79$ already at $p = 16$ --- while a corrected profile, the geometric climb with a fixed additive offset $1/(\LL - 1) = 1.70951$ per block, satisfies all $p$ size conditions by construction, with no correction step. The two halves compose to a proof for every $p \ge 16$, at $\gamma$ between the absolute constants $3.683012$ and $5.140212$ --- no $p$-dependence at all --- with $3 \le p \le 15$ meeting the same bracket by finite check and $p \in \{2,4\}$, outside the construction's reach, covered by direct exhibition. Every constructed configuration remains a size-passer only and fails the divisibility conditions $q \mid R_r$, as the instances above do: sharper evidence that counting cannot do better, and no evidence about exclusion. The record is \texttt{cycles.md} \S12.8.6 (\url{https://github.com/macindoe/collatz/blob/9d9d1ec/cycles.md}); the note added in v2 and the continued-fraction route it named are in the release description and at \S12.8.6.3 there.

\section{The equidistribution hypothesis, made exact}\label{sec:aeh}

What is the right comparison object? By Theorem~\ref{thm:onestep} every dynamical decision at horizon one is a function of the depth-$k$ window, so the null model is a measure on windows, and the structure forces the choice: the $\w$-residues carry the anchor digits, for which Haar measure is the canonical reference (this is where the classical $2$-adic shift-conjugacy heuristics live), while the depth is \emph{dynamical} --- it is produced by Theorem~\ref{thm:depth} --- and receives instead its own exact renewal law, $\dnext = m_+ + a_+$ with $m_+$ geometric$(1/2)$ and \emph{independent} of the absorption $a_+$, so that it is the distribution of the depth, and not that of the absorption, which is the convolution of the two. The stationary $3$-adic law governing $a_+$ (\texttt{aeh.md} \S13.6.5) is not new here: it is Tao's Syracuse random variable $\mathrm{Syrac}(\Z_3)$ \cite[Lemma~1.12 and Remark~1.13]{tao}, in the present normalisation. Precisely, the $3$-adic past-limit of the block coordinate has the law of $\mathrm{Syrac}(\Z_3)/2$ --- the two differ by the $3$-adic unit $2^{-1}$ because one block here is a run of Syracuse steps terminated by the first exponent $\ge 2$, so that the block chain sees the Syracuse chain conditioned on that event --- and hence $a_+ = \vth{\mathrm{Syrac}(\Z_3) + 2}$, with $P(a_+ = 0) = \tfrac23$, $P(a_+ = 1) = \tfrac{19}{63}$ and $P(a_+ \ge 2) = \tfrac{2}{63}$ read off the nine values of $\mathrm{Syrac}(\Z/9\Z)$ that Tao computes. What the present coordinates add is the \emph{joint labelled} law (\texttt{aeh.md} \S13.6.3(v)): the $\w$-residues are Haar-uniform among odd residues and independent of the depth, the whole carried on the residues \emph{together with} the stratum labels $(s,\sigma,a_+)$ --- $a_+$ being a $3$-adic function of $\w$ (Lemma~\ref{lem:absorption}) which the finite $2$-adic residue window does not determine. Two integers fix the observable: a depth $k$ and a \emph{cap} $D$, chosen together and quantified over. The \emph{capped depth-$k$ window} at a visit is
\[
  W_{k,D} \;=\; \bigl(\, \w \bmod 2^{k+2},\; d \wedge D,\; s \wedge D,\;
  \sigma \wedge D,\; a_+ \wedge D \,\bigr), \qquad u \wedge D := \min(u,D),
\]
a finite alphabet of at most $2^{k+1}D^3(D+1)$ letters; the cap is what bounds it, and it must cap the labels as well as the depth, since $\sigma$ and $a_+$ are unbounded. It is not the window of Theorem~\ref{thm:onestep}, whose alphabet is countably infinite and which decides the next step: $W_{k,D}$ decides nothing, and no statement in this section asks it to. Let $\pi_{k,D}$ denote the law of $W_{k,D}$ under the two-sided Bernoulli measure $\hat B = \bigotimes_{i \in \Z}(\mathrm{geom}(\tfrac12) \times \mathrm{geom}(\tfrac12))$ on the door-letter alphabet --- two-sided because the absorption $a_+$ is a function of the orbit's $3$-adic past and has no realisation on the one-sided $2$-adic side (\texttt{aeh.md} \S13.6.3(iii), \S13.6.5); reading the stationary Syracuse variable as the outcome of an iteration extending to arbitrarily large negative times is Tao's own alternative to his positive-time indexing \cite[Remark~1.13, footnote~4 of arXiv~v7]{tao}. ``Product'' names exactly two clauses of $\pi_{k,D}$ --- residue $\perp$ depth, and $m_+ \perp a_+$ --- and no others: $\pi_{k,D}$ is not a product across its coordinates ($s$ is an exact function of the full state), and the process $(W_{k,D}(n))_n$ is stationary but not independent across time (e.g.\ from $(\w \equiv 1\ (8),\ d$ odd$)$ the next depth is exactly $1$). Under $\pi_{k,D}$, unconditionally: $P(s = j) = 2^{-j}$ at every $j < D$, the mass beyond the cap sitting in one tail cell; the $3$-gain rate is $\sum_{j\,\mathrm{even}} 2^{-j} = \tfrac13$ by Lemma~\ref{lem:absorption}, exact up to that same cell, whose parity split the cap hides; and the classical negative drift is a property of $\pi_{k,D}$ --- but only of $\pi_{k,D}$: window equidistribution at each fixed $(k,D)$ does not control the means of the unbounded $m_+$ and $s$, so no drift or contraction statement about orbits follows from it (\texttt{aeh.md} \S13.3.2). The folklore ``ledger'' is thus a theorem \emph{about} $\pi_{k,D}$; the empirical question is only whether the orbits of typical starting values follow $\pi_{k,D}$.

\begin{hypothesis}[AEH, ensemble form]\label{hyp:aeh}
Fix a \emph{budget rate} $\tau > 0$ and a \emph{block horizon rate} $\theta > 0$;
for each $N$ put $b = \lfloor\log_2 N\rfloor$, $\Lambda_N = \lceil \tau b\rceil$
and $T_N = \lceil \theta b\rceil$. Draw $x$ uniformly from the odd integers of
$[N, 2N)$; put $(\w_0,d_0) = R(x)$ and $(\w_{n+1},d_{n+1}) = F(\w_n,d_n)$; let the
\emph{letter} at block $n$ be $\ell_n = (m_{+,n},\,s_{n+1})$, and let
$S_n = \sum_{i<n}(m_i + s_i)$ be the \emph{exponent spent} before block $n$ ---
the number of $2$'s divided out from $x$ to $x_{\mathrm{exit}}(n-1)$. Tally block
$n$ at the symbol
\[
  \tilde\ell_n \;=\; \ell_n \ \text{ if } S_n < \Lambda_N, \qquad
  \tilde\ell_n \;=\; \dagger \ \text{ otherwise.}
\]
For a finite word $w = (w_1,\dots,w_L)$ of letters, let $f_N(w,x)$ be the
frequency of $w$ among the $T_N - L + 1$ blocks
$(\tilde\ell_n,\dots,\tilde\ell_{n+L-1})$, $0 \le n \le T_N - L$: every
block counted exactly once, at weight $1/(T_N-L+1)$ --- the same number for
every block of every orbit, so no block is reweighted by the orbit it came from
and no denominator is random. Then for every finite word $w$ and every
$\varepsilon > 0$,
\[
  \frac{2}{N}\,\#\bigl\{\, x \text{ odd},\ N \le x < 2N \;:\;
  \bigl| f_N(w,x) - B[w] \bigr| > \varepsilon \,\bigr\}
  \;\longrightarrow\; 0 \qquad (N \to \infty),
\]
where, writing $w_i = (m_i, r_i)$ for the components of the $i$-th letter of $w$,
$B[w] = \prod_{i} 2^{-(m_i + r_i)}$, and $B[w] = 0$ for any $w$ containing
$\dagger$, for every admissible pair $(\tau,\theta)$ in the sense defined below.
\end{hypothesis}

Equivalently, and this is the form the calibration measures: for every $k$, $D$
and $L$, the empirical distribution of the $L$-blocks of consecutive capped
windows $W_{k,D}$ over the same tallied blocks, with $\dagger$ carried through
and given mass $0$, converges in total variation on the finite window alphabet,
off a vanishing density of starts, to $\pi^{(L)}_{k,D}$, the $L$-block law of the
stationary process under $\hat B$.
The equivalence is Theorem 13.6.4 of \texttt{aeh.md}, a deterministic dictionary
between letters and labelled window blocks; the case $L = 1$ recovers
$\pi_{k,D}$, and $\pi^{(L)}_{k,D} \neq \pi_{k,D}^{\otimes L}$. A sampled segment
has no infinite past at its first block, so the reconstruction of
$W_{k,D}$ from letters is exact only away from the start; the
deviation occupies $O(1)$ blocks of a horizon $T_N \to \infty$ and contributes
$O(1/T_N)$ to every frequency above.

There is one limit here, $N \to \infty$, and the sample grows because the sampling
scale grows rather than because any one orbit is run forever. That is forced: along
a fixed orbit the unrestricted empirical distribution is false, every convergent
tail sitting at $(1,1)$ forever, while above a fixed cut a convergent orbit
supplies only finitely many qualifying visits and none at all once the cut exceeds
its maximum, so a limit in orbit length is empty rather than merely delicate, in
whichever order it is taken (\texttt{aeh.md} \S13.6.6).

The horizon does the job a bulk cut would do, and does it without one. A step of
the one-division map $y \mapsto y/2$ or $(3y+1)/2$ never lowers $\log_2$ by more
than $1$, so $\log_2 x_{\mathrm{exit}}(n-1) \ge \log_2 x - S_n$ for every start and
every $n$, with no exceptional set; inside a budget $\tau < 1$ of the start's own
bits every tallied exit therefore exceeds $N^{1-\tau}$, far above the
\emph{bottom regime} of \texttt{aeh.md} \S13.1. That matters because any altitude
threshold is a selection on the observable itself --- a cut on
$x_{\mathrm{exit}} = (3^{d}\w - 1)/2^{s}$ censors large $s$, and a cut on the core
$\wnext$ censors large $s$, $m_+$ and $a_+$ together --- whereas the budget clause
$S_n < \Lambda_N$ is \emph{predictable}, decided by blocks strictly earlier than
the one it admits, and the tally denominator is the deterministic $T_N$ that
\texttt{aeh.md} \S13.5's standing rule was written to secure. Call $\tau$
\emph{protected} when all but a
vanishing density of starts keep every in-budget exit above any $X_N$ with
$\log X_N = o(\log N)$, and $(\tau,\theta)$ \emph{consistent} when
$\theta\,\mathbb{E}_B[m+r] < \tau$; \emph{admissible} means both. Every
$\tau < 1$ is protected outright by the bound above, and every
$\tau < (1-\log_2\sqrt3)^{-1} = 4.8188\ldots$ is protected at natural density one
by Inselmann \cite[Thm.~1.1]{inselmann}, whose two-sided envelope is stated for
the one-division map --- the same unit, nothing converted.
Consistency is compatibility with the target law rather than a claim about
orbits: under $B$ a block spends $\mathbb{E}_B[m+r] = 4$ of exponent, so
$4\theta < \tau$ says exactly that $B$ itself predicts no out-of-budget block,
which is what $\pi_{k,D}$'s giving $\dagger$ no mass requires. Where the cylinder
count runs it is more, and unconditionally: for $\tau < 1$ no block of the horizon
leaves the budget and the empirical exponent mean converges to $4$
(\texttt{aeh.md} Lemma 13.2.4(g), whose two error terms are, on that range,
precisely the two clauses of admissibility). Past that range the hypothesis
supplies less, and only about frequencies: a vanishing frequency of $\dagger$ says
all but $o(T_N)$ of the first $T_N$ blocks are within budget, and is not a bound on
a sum over its complement --- the block at which the budget is exhausted is tallied
at full weight and its letter is unbounded, and the $o(T_N)$ blocks past it are
uncontrolled, so $T_N^{-1}\sum_{n<T_N}(m_n+r_n) \to 4$ does not follow there and
can fail by any amount. Hence $4\theta < \tau < 1$ is $\theta < 1/4$ and
$4\theta < \tau < 4.8188\ldots$ is $\theta < 1/\beta = 1.2047\ldots$, where
$\beta = 2(2-\LL) = 0.8301\ldots$ is the classical per-block contraction rate and
$4/\beta = (1-\log_2\sqrt3)^{-1}$ an identity --- arithmetic on the definition of
consistency, with the divisor in either block reading the mean of the target law:
a theorem about orbits below the budget, and not a consequence of the hypothesis
above it.

The hypothesis has an unconditional base case, and it is Heuristic~\ref{prop:budget}
with a number on it. The classical coding fact --- the odd integers whose first $n$
blocks realize a prescribed itinerary form exactly one residue class modulo
$2^{S+1}$, $S$ the itinerary's total exponent (Terras \cite{terras}; in the present
coordinates \texttt{itinerary.md} \S14.15.1.5) --- makes the first $n$ blocks of a
start drawn uniformly from \emph{any} window of $2^{J}$ consecutive integers
exactly product-distributed on the event $S + 1 \le J$: the event that the itinerary
has not outspent the start's supply of binary digits. That fact is about the word
beginning at $x$ itself, while $\ell_0$ is the letter of the block \emph{after} the
start's own; it is applied here to the \emph{extended} $(n+1)$-letter word --- the
start's own letter followed by $\ell_0,\dots,\ell_{n-1}$, of total exponent exactly
$S_{n+1}$, the divisions from $x$ down to $x_{\mathrm{exit}}(n)$ --- of whose law
the one displayed is a marginal. Writing
$b = \lfloor\log_2 N\rfloor$, this gives, for every $J \ge 2$,
\[
  \begin{gathered}
    \bigl\lVert \mathrm{Law}(\ell_0,\dots,\ell_{n-1}) - B^{\otimes n}
      \bigr\rVert_{\mathrm{TV}}
    \;\le\; \frac{2^{J+2}}{N} + P_B(S_{n+1} \ge J), \\[2pt]
    P_B(S_{n+1} \ge J) = P\bigl(\mathrm{Bin}(J-1,\tfrac12) < 2(n+1)\bigr),
  \end{gathered}
\]
the second identity because $S_{n+1}$ is the waiting time for the $2(n+1)$-th
head in a fair coin sequence. The first term is the price of a general window
$[N,2N)$ over a dyadic one and is negligible for $J \le (1-\eta)b$; the second
decays at the exact rate $\log 2 - H(2\theta)$ per bit at $n = \lceil\theta b\rceil$,
$H$ the binary entropy --- positive for every $\theta < 1/4$ and \emph{zero} at
$\theta = 1/4$ (\texttt{aeh.md} \S13.2). The bound is on the joint law of the whole
length-$n$ word, so it controls every finite block length at once, and with the
concentration of empirical pattern frequencies, the removal of the segment's initial
past-boundary and the altitude bound above --- which makes the budget its own
protection --- it yields the hypothesis outright:
Hypothesis~\ref{hyp:aeh} is a \emph{theorem} for every admissible $(\tau,\theta)$
with $\tau < 1$, equivalently for every horizon rate $\theta < 1/4$, at every
block length. The assembly is Lemma 13.2.4 of \texttt{aeh.md}.

The natural clock here is total exponent, not blocks: the frontier is one unit of
exponent per bit of start --- the start's own digit supply, with no estimate in it
--- and it reads as $1/4$ blocks per bit only through the mean exponent per block,
$\mathbb{E}_B[m+r] = 4$, a conversion available here precisely because the word is
exactly $B$ here. Hypothesis~\ref{hyp:aeh} is the assertion that the law still
describes the orbit after the start's digits have been spent; the digit budget
locates the frontier, and this is the statistical statement on the far side of it.
That frontier is the classical one --- the same $\theta < 1/4$ is where Terras's
count stops, and $1 - \beta/4 = \log_4 3$ is exactly the exponent Korec
\cite{korec} obtains by exhausting it --- and it is a frontier of this technique
and not of the problem, Inselmann \cite{inselmann} crossing it unconditionally for
the trajectory envelope and the first moment, in step time and by a stronger
density notion rather than a sharper count.

AEH implies the ledger, in the form the hypothesis has: $s$ is a coordinate of
the observable, so $P(s = j) = 2^{-j}$ is a marginal of $\pi_{k,D}$ rather than a
deduction, exactly at every $j < D$ and up to the cap's single tail cell above;
and for every $\varepsilon$ and every horizon rate, all but a vanishing density of
the odd starting values of each size carry those frequencies along their first
$T_N$ blocks, all but $o(T_N)$ of them within the budget --- and all of them when
$\tau < 1$. Evaluating
each odd $x$ once, at its own dyadic scale $2^{\lfloor \log_2 x\rfloor}$, turns
that array of statements into a single one: the odd integers that fail at their
own scale have natural density zero (\texttt{aeh.md} Proposition 13.2.5). The
union over \emph{all} sampling scales is a different and larger object and is not
controlled --- the bad sets form a triangular array, each testing a horizon and a
budget that move with the scale, and vanishing density in every window does not
thin their unrestricted union. The
descent consequence is not stated here, because it is a theorem without the
hypothesis and a stronger one: Inselmann's \cite[Cor.~1.4, Thm.~1.10]{inselmann},
cited in Related work, is two-sided, uniform in the time, unconditional, and runs
to $4.8188\ldots$ times the classical range at which this section's cylinder count
stops, the two measured in the same units. His \cite[Thm.~1.6]{inselmann} is what
carries his own passage from $T$-time to Syracuse time; the further passage to
block time needs the frequency with which a Syracuse step ends a block, a
two-letter statistic of the parity word that neither theorem supplies and that
$\pi_{k,D}$ does, so it is not available to underwrite the hypothesis: inside the
digit budget it is a theorem of the cylinder count (\texttt{aeh.md}
Lemma 13.2.4(g)), and past it neither a theorem nor a consequence of the hypothesis
(\texttt{aeh.md} \S13.2.3). What the hypothesis supplies beyond these is
distributional and is the whole of its content: the full $2^{-j}$ marginal at every
$j$ below the cap, the exact $\tfrac13$ rate of Lemma~\ref{lem:absorption}, and the
depth law, at every finite depth and past the budget at which the classical count
stops. Its exceptional set is a set of \emph{starting values} of vanishing density
at a prescribed finite horizon --- not a null set of orbits, and not a statement
about any orbit's infinite tail. It does not exclude individual staircase tails
(Remark~\ref{rem:staircase}); it does not iterate --- the image of a density-one set
of starts need not be density-one at the next scale, an obstruction Tao states in the
same terms \cite{tao}; and by Heuristic~\ref{prop:budget} it cannot be reached by
finite-window computation. Its content is a question about digits of $2$-adic
logarithms, older and broader than the Collatz problem.

\paragraph{Calibration.} An eight-round campaign (per-orbit statistics over each cell's own conditional counts, $1{,}600$--$2{,}600$ independent orbits per cell, under an altitude cut on the core) produced four apparent violations of $\pi_{k,D}$-uniformity, each dissolved by a sharper control: an $s{=}3$ marginal excess and pair correlations (artifacts of small starting values); digit biases (contamination by the \emph{bottom regime}, the fixed drainage basin of small integers, where deviations reach $z=41$ but reflect particular numbers rather than a measure); and a $5\sigma$ digit bias in one cell surviving a $2^{40}$ cut. The last dissolved under fixed-horizon unweighted sampling ($0.2503\pm0.0015$ against null $0.25$ over $84{,}739$ visits) and was explained exactly:

\begin{lemma}[deterministic routing]\label{lem:routing}
For $\w \equiv 1 \pmod 8$, $d = 1$: here $s = 1$, $C = 3\w + 1$, $\vt{C} = 2$, $a_+ = 0$, so $\wnext = (3\w+1)/4$ and $\dnext = 1$; the residue $\w \bmod 32$ determines the next class exactly: $1 \mapsto (1,1)$ (immediate self-return), $9 \mapsto (7,1)$, $17 \mapsto (5,1)$, $25 \mapsto (3,1)$, the lifting class.
\end{lemma}

Ratio estimators over correlated visit sequences inherit this routing (residue $25$ leads into deep descent, hence fewer qualifying visits, hence inflated per-orbit weight), manufacturing a bias aimed at one residue of the unique class with a self-loop channel. No residual discrepancy was detected at any tested depth or cell, under the stated protocol and within three limits the campaign does not reach past: it tests block lengths $L = 1$ and $L = 2$ only and is silent above; its adjudicating runs pool per visit across orbits, so they measure the across-orbit average of the per-start frequencies --- a consequence of Hypothesis~\ref{hyp:aeh} rather than the per-start statement itself, which would be tested by the distribution across orbits of $\lVert\nu - \pi_{k,D}\rVert_{\mathrm{TV}}$ and is not among the runs reported; and its altitude guard cuts on the core $\wnext$ rather than the door, censoring $s$, $m_+$ and $a_+$ at once, and does bind at finite size, removing $2.6\%$ of visits in the adjudicating run and moving the reported statistics at the third decimal (\texttt{aeh.md} \S13.4, \S13.5). Hypothesis~\ref{hyp:aeh} carries no cut of any kind; it tallies within the exponent budget, which bounds every tallied exit from below without selecting on the letter being tallied.

\section{Discussion}

In reduced coordinates the Collatz problem factors cleanly. The local arithmetic is under exact control: one $2$-adic logarithm governs the step, and finite windows decide steps without error. What remains is a single question wearing two costumes: for typical starting values, whether the anchors' digits equidistribute (Hypothesis~\ref{hyp:aeh} --- statistics); for cycles, whether closed anchor walks are rigid (Theorem~\ref{thm:staircase} --- arithmetic). The staircase shows the costumes cover the same body. Three directions have defined success criteria: rigidity statements for closed anchor walks beyond the size level (the only route past Theorem~\ref{thm:staircase}); transfer-operator analysis of the exact window chain; and any progress on digits of $2$-adic logarithms of integers, a problem this paper inherits rather than creates.

\appendix
\section{Methodology, responsibility, and records}\label{app:method}
This work was developed through extensive AI-assisted research over three years (2023--2026), from a hand-drawn origin note (13 March 2023) through seventy-nine drafts to a research wiki whose pages carry status headers, verification records, and refuted claims with their refutation data. The author is responsible for the published claims, the research direction, and the standards of evidence; the AI collaborator carried out the mathematical development within the author's coordinate system --- formulations, proofs, computations, and connections to the literature. Because this division of labor is unusual, correctness claims rest on the documented verification protocol rather than on authorial authority alone: every computational claim cites a runnable script; every result labeled proved was re-verified with independently written code before being so labeled; the calibration campaign's four dissolved anomalies, and two corrected search-filter errors on the cycle side, are documented rather than smoothed over. The complete record (wiki, code, draft history, origin note) is public at \url{https://github.com/macindoe/collatz}; every wiki section and script named in this paper is cited at commit \texttt{1663d30}.

\begin{thebibliography}{99}
\small
\bibitem{terras} Riho Terras, \emph{A stopping time problem on the positive integers}, Acta Arith.\ 30 (1976), 241--252.
\bibitem{everett} C.~J.~Everett, \emph{Iteration of the number-theoretic function
  $f(2n) = n$, $f(2n+1) = 3n+2$}, Adv.\ Math.\ 25 (1977), 42--45.
\bibitem{korec} I.~Korec, \emph{A density estimate for the $3x+1$ problem},
  Math.\ Slovaca 44 (1994), no.\ 1, 85--89.
\bibitem{tao} T.~Tao, \emph{Almost all orbits of the Collatz map attain almost
  bounded values}, Forum Math.\ Pi 10 (2022), Paper No.\ e12, 56 pp.;
  arXiv:1909.03562.
\bibitem{inselmann} M.~Inselmann, \emph{An approximation of the Collatz map and a
  lower bound for the average total stopping time}, arXiv:2402.03276 (2024).
\bibitem{wirsching} G.~J.~Wirsching, \emph{The Dynamical System Generated by the
  $3n+1$ Function}, Lecture Notes in Math.\ 1681, Springer-Verlag, Berlin, 1998.
\bibitem{thomas} A.~Thomas, \emph{A non-uniform distribution property of most
  orbits, in case the $3x+1$ conjecture is true}, Acta Arith.\ \textbf{178}
  (2017), no.~2, 125--134; arXiv:1512.05852.
\bibitem{steiner} R.~P.~Steiner, \emph{A theorem on the Syracuse problem}, Proc.\ 7th Manitoba Conf.\ Numerical Math.\ (1977), Utilitas Math., Winnipeg, 1978, pp.\ 553--559.
\bibitem{eliahou} S.~Eliahou, \emph{The 3x+1 problem: new lower bounds on nontrivial cycle lengths}, Discrete Math.\ 118 (1993), 45--56.
\bibitem{sdw} J.~Simons, B.~de~Weger, \emph{Theoretical and computational bounds for $m$-cycles of the $3n+1$ problem}, Acta Arith.\ 117 (2005), 51--70.
\bibitem{hercher} C.~Hercher, \emph{There are no Collatz $m$-cycles with $m \le 91$}, J.\ Integer Seq.\ 26 (2023), Article 23.3.5; arXiv:2201.00406.
\bibitem{lagarias} J.~C.~Lagarias, \emph{The 3x+1 problem: an annotated bibliography}, I--II, arXiv math/0309224, math/0608208.
\bibitem{yu} K.~Yu, \emph{Linear forms in $p$-adic logarithms}, Acta Arith.\ 53 (1989), 107--186; \emph{Linear forms in $p$-adic logarithms II}, Compositio Math.\ 74 (1990), 15--113; \emph{Linear forms in $p$-adic logarithms III}, Compositio Math.\ 91 (1994), 241--276.
\bibitem{bl} Y.~Bugeaud, M.~Laurent, \emph{Minoration effective de la distance $p$-adique entre puissances de nombres alg\'ebriques}, J.\ Number Theory 61 (1996), 311--342.
\bibitem{rhin} G.~Rhin, \emph{Approximants de Pad\'e et mesures effectives d'irrationalit\'e}, in: S\'eminaire de Th\'eorie des Nombres, Paris 1985--86, Progr.\ Math.\ 71, Birkh\"auser, Boston, 1987, pp.\ 155--164 (the Proposition used here is on p.~160); see also Q.~Wu, \emph{On the linear independence measure of logarithms of rational numbers}, Math.\ Comp.\ 72 (2003), no.\ 242, 901--911.
\bibitem{barina} D.~Barina, \emph{Improved verification limit for the convergence of the Collatz conjecture}, J.\ Supercomputing 81 (2025), article 810; DOI 10.1007/s11227-025-07337-0.
\bibitem{merle} E.~Merle, \emph{On the non-existence of non-trivial Collatz cycles: a conditional formal proof in Lean 4 with documented structural obstructions}, preprint, Zenodo DOI 10.5281/zenodo.19790406 (2026).
\bibitem{llmcollatz} E.~Y.~Chang, \emph{Exploring Collatz Dynamics with Human--LLM Collaboration}, arXiv:2603.11066 (2026).
\end{thebibliography}

\end{document}
